\documentclass[a4paper,12pt]{article}

% --- Use XeLaTeX font support ---

\usepackage{fontspec}
\setmainfont{TeX Gyre Termes} % nice Times-like font, installed by texlive-core
\setsansfont{TeX Gyre Heros} % optional sans-serif
\setmonofont{Fira Code}       % optional monospaced font
% --- Math packages ---
\usepackage{amsmath, amssymb, amsthm}
\usepackage{mathtools}

% --- Page layout ---
\usepackage[margin=1in]{geometry}

% --- Better lists ---
\usepackage{enumitem}
% --- Hyperlinks for references ---
\usepackage[colorlinks=true, linkcolor=blue, urlcolor=blue]{hyperref}

% Header: fancy layout like the rovided image
\usepackage{fancyhdr}
\setlength{\headheight}{26pt} % increase if fancyhdr warns
\pagestyle{fancy}
\fancyhf{} % clear header/footer
\rhead{\begin{tabular}{r} Mt.: 3052737\end{tabular}}
\chead{\begin{tabular}{c}Lösungen EA 2\end{tabular}}
\lhead{\begin{tabular}{l}61521 Einführung in die \\ Numerische Mathematik SS26 \end{tabular}}
\renewcommand{\headrulewidth}{0.8pt}
\fancyfoot[C]{- \thepage\ -}
% Ensure the plain page style (used by \maketitle) uses the same header

% --- Line spacing ---
\usepackage{setspace} % provides \doublespacing and \onehalfspacing

% --- Optional solutions environment ---
\newenvironment{solution}{
  \par\noindent\textbf{Solution:}\quad
}{\hfill$\blacksquare$\par\vspace{1em}}

% % --- Title info ---
% \title{Aufgabe \#1}
% \author{Your Name}
% \date{\today}
% 
\begin{document}
\raggedright
% Enable double (2.0) line spacing for the document body
\doublespacing
\noindent\textbf{Aufgabe 1} \\
\noindent\textbf{(a)\\(i)} Gegeben $\varphi'(x^*) = 0$. Aufgrund der Stetigkeit von $\varphi'$ existiert für jedes $\kappa \in (0, 1)$ eine Umgebung $U_\delta = (x^* - \delta, x^* + \delta)$ mit $\delta > 0$, sodass für alle $x \in U_\delta$ gilt:$$|\varphi'(x)| \le \kappa$$
Für zwei beliebige Punkte $x, y \in U_\delta$ existiert nach dem Mittelwertsatz der Differentialrechnung ein $\xi$ zwischen $x$ und $y$, für das gilt:$$\varphi(x) - \varphi(y) = \varphi'(\xi)(x - y)$$
Da  $\xi$ in $U_\delta$ liegt, folgt $|\varphi'(\xi)| \le \kappa$. Daraus:
$$|\varphi(x) - \varphi(y)| = |\varphi'(\xi)| \cdot |x - y| \le \kappa |x - y|.$$
Also ist $\varphi$ auf $U_\delta$ kontrahierend mit der Konstanten $\kappa \in (0, 1)$.
Für beliebiges $x \in U_{\delta}$ gilt
$$|\varphi(x) - x^*| = |\varphi(x) - \varphi(x^*)| \le \kappa |x - x^*| < \kappa \delta < \delta$$
 Somit liegt $\varphi(x)$ wieder in $U_\delta$. $\varphi: U_\delta \to U_\delta$ ist demnach eine kontrahierende Selbstabbildung. \\
 \bigskip
 
 Da $\varphi$ auf der Umgebung $U_\delta$ eine kontrahierende Selbstabbildung ist, liefert der Banachsche Fixpunktsatz, dass die Iterationsfolge $x_{k+1} = \varphi(x_k)$ für jeden Startwert aus dieser Umgebung wohldefiniert ist und gegen den eindeutigen Fixpunkt $x^*$ konvergiert. Es folgt die lokale Konvergenz des Verfahrens. \\
 \bigskip
 \noindent\textbf{(ii)} 
\begin{itemize}
  \item Da $\varphi$ $(p+1)$-mal stetig differenzierbar ist, existiert nach dem Satz von Taylor ein $\xi$ zwischen $x$ und $x^*$, sodass:$$\varphi(x) = \varphi(x^*) + \sum_{j=1}^{p} \frac{\varphi^{(j)}(x^*)}{j!} (x - x^*)^j + \frac{\varphi^{(p+1)}(\xi)}{(p+1)!} (x - x^*)^{p+1}$$
  Bildung des Betrags:
  $$|\varphi(x) - \varphi(x^*)| = \frac{|\varphi^{(p+1)}(\xi)|}{(p+1)!} |x - x^*|^{p+1}$$
  Einsetzen der Ungleichung aus der Aufgabenstellung ($|\varphi^{(p+1)}(x)| \le M(p+1)!$):
  $$\boxed{|\varphi(x) - \varphi(x^*)| \le \frac{M(p+1)!}{(p+1)!} |x - x^*|^{p+1} = M |x - x^*|^{p+1}}$$

  \item Einsetzen von   $x_{k+1} = \varphi(x_k)$ in die Ungleichung:
$$|x_{k+1} - x^*| \le M |x_k - x^*|^{p+1}$$
$$\frac{|x_{k+1} - x^*|}{|x_k - x^*|^{p+1}} \le M.$$
Also
$$\limsup_{k \to \infty} \frac{|x_{k+1} - x^*|}{|x_k - x^*|^{p+1}} \in \mathbb{R}.$$
Das beweist, dass die Q-Ordnung mindestens $p+1$ ist.
\end{itemize}

\bigskip
\noindent\textbf{(b)}
Das Newton-Verfahren hat die Iterationsfunktion:$$\varphi(x) = x - \frac{f(x)}{f'(x)}$$
Die erste Ableitung ist:
$$\varphi'(x) = 1 - \frac{f'(x)f'(x) - f(x)f''(x)}{[f'(x)]^2} = \frac{f(x)f''(x)}{[f'(x)]^2}$$
Einsetzen von $f(x^*) = 0$:
$$\varphi'(x^*) = \frac{0 \cdot f''(x^*)}{[f'(x^*)]^2} = 0$$
Da $f$ dreimal stetig differenzierbar ist, ist $\varphi$ in einer Umgebung von $x^*$ zweimal stetig differenzierbar. \\ 
Die Voraussetzungen aus Aufgabenteil (a) für $p=1$ sind erfüllt. Nach (a)(ii) ist die Q-Ordnung des Verfahrens somit mindestens $p+1 = 1+1 = 2$. Das Verfahren konvergiert lokal Q-quadratisch. \\

\bigskip
\noindent\textbf{(c)}
Der erste Schritt ($y_n$) ist das Standard-Newton-Verfahren. Aus (b) wissen wir, dass dessen Fehler quadratisch abfällt, $\tilde{e}_n = \mathcal{O}(e_n^2)$.
Für den Fehler im zweiten Schritt betrachtet man $e_{n+1} = x_{n+1} - x^*$:$$e_{n+1} = y_n - x^* - \frac{f(y_n)}{f'(x_n)} = \tilde{e}_n - \frac{f(y_n)}{f'(x_n)} = \frac{\tilde{e}_n f'(x_n) - f(y_n)}{f'(x_n)}$$
Taylor-Reihen um $x^*$: \\
Für $f'(x_n)$:\\
$$f'(x_n) = f'(x^*) + f''(x^*)e_n + \mathcal{O}(e_n^2)$$
Für $f(y_n)$:\\
Wegen $f(x^*) = 0$ gilt:
$$f(y_n) = f'(x^*)\tilde{e}_n + \mathcal{O}(\tilde{e}_n^2)$$
Da $\tilde{e}_n = \mathcal{O}(e_n^2)$ ist, wird $\mathcal{O}(\tilde{e}_n^2) = \mathcal{O}(e_n^4)$. Somit:
$$f(y_n) = f'(x^*)\tilde{e}_n + \mathcal{O}(e_n^4)$$
Einsetzen in den Zähler:
\begin{align*}
\text{Zähler} &= \tilde{e}_n \left[ f'(x^*) + f''(x^*)e_n + \mathcal{O}(e_n^2) \right] - \left[ f'(x^*)\tilde{e}_n + \mathcal{O}(e_n^4) \right] \\
&= \tilde{e}_n f''(x^*)e_n + \tilde{e}_n \mathcal{O}(e_n^2) + \mathcal{O}(e_n^4) \\
 &= \mathcal{O}(e_n^2)\,e_n + \mathcal{O}(e_n^2)\mathcal{O}(e_n^2) + \mathcal{O}(e_n^4) \\
 &= \mathcal{O}(e_n^3)
\end{align*}
Der Nenner $f'(x_n)$ nähert sich $f'(x^*) \neq 0$ an, also einer Konstanten. Daher ergibt sich für den Gesamtfehler:$$e_{n+1} = \frac{\mathcal{O}(e_n^3)}{f'(x^*) + \mathcal{O}(e_n)} = \mathcal{O}(e_n^3)$$
Damit ist gezeigt, dass das zweistufige Newton-Verfahren lokal mindestens die Q-Ordnung 3 besitzt.
\end{document}
